\section{章节罗盘：扩张的能力与扩张的账目}
\label{sec:high-qing-compass}

清初的征服建立了多区域帝国的框架，本章考察雍正、乾隆时期如何使这一框架运转得更快、更远。财政整顿、密折与军机文书、驻防和地方督抚构成一条协调链，却不能取消距离、季节、道路和地方合作者对政策的限制。所谓“盛世”在这里首先是一种能够调兵、核销、赈济和经营边地的能力，而非没有摩擦的稳定状态。

读者将随这条链进入西北、青藏、金川与海疆，并回到河工、盐政、漕运和身份供养的日常账册。军事胜利、设治和社会秩序须分开衡量：每一次把边地纳入帝国的行动，都增加夫役、军费、移民安置与驻防供给。本章的收束将说明，扩张没有消除清初留下的治理成本，而是扩大了它们的地理半径。
